\guidechapter{Putting exactness and the running time together}{Read before the final assembly proof.}

The paper has two finish lines. It must return the exact product, and every
part of the computation must fit the claimed time bound. These obligations
are related through precision and storage volume, but neither replaces the other.

\paragraph{Exactness from a strict error interval.}
The desired convolution coefficients $c_j$ are integers. If the computation
produces an approximation $\widetilde c_j$ with
$|\widetilde c_j-c_j|<1/2$, nearest-integer rounding recovers $c_j$ uniquely.
The bound must hold \emph{after} restoring all transform normalizations and
digit scales. A tie at exactly $1/2$ is not enough. Once coefficients are
recovered, carries and binary output are exact integer operations.

\explfigure{09-rounding}{Approximate internal arithmetic, exact final output.}{
The displayed $12.27$ is a teaching example. The source's proof supplies a
uniform bound on every recovered coefficient; this picture does not measure
its actual numerical error.}

The source uses radix width $b=\lceil\log_2 n\rceil$ and working precision
$p=6b$ in its original assembly. Roughly, there are $n/b$ digits and
$\Theta(p)$ stored bits per coefficient, making the padded volume
$V=Tp=\Theta(n)$. This relation is what turns costs per stored bit into a
bound in terms of the original input size.

\paragraph{Seven competing bills.}
The final table separates operations that scale differently with $p$.
Write a normalized term as $p^{1-g_i}$; then $g_i$ is its margin below a
full power of $p$. With $d=\Theta(p^\epsilon)$,
$K=\Theta(p^{\epsilon c})$, and $\ell=\Theta(p^{1-\epsilon})$,
the margins are
\begin{align*}
 g_1&=1-\epsilon(1+c),&g_2&=\epsilon c(1-\tau),\\
 g_3&=\epsilon(1-\lambda'),&g_4&=1-\tau-\epsilon(2-\tau),\\
 g_5&=\tfrac14-\delta-\tfrac32\epsilon,&g_6&=1-\delta-\epsilon,\\
 g_7&=\epsilon.&&
\end{align*}
The total is controlled by the \emph{smallest} margin. Making an already
large margin larger may have no effect on the headline exponent.

\explfigure{10-margins}{The original parameter choice has a clear bottleneck.}{
These values are computed from exact rational parameters in the original
source, not estimated from timings. The horizontal scale is $-\log_2g_i$:
longer bars mean smaller savings. Chunk exchanges have $g_2=2^{-181}$;
the chosen final saving is $\kappa=2^{-182}$.}

\paragraph{Read one row all the way through.}
For transform-layout exchanges, there are $O(p/K)$ swaps of width $K$.
Each costs $O(VK^\tau)$. Thus the total is
\[
 O(VpK^{\tau-1})
 =O\bigl(Vp^{1-\epsilon c(1-\tau)}\bigr).
\]
This is exactly the second margin $g_2$. The small quantities multiply:
a small width-growth exponent $\epsilon c$ and a small swap saving
$1-\tau$ yield an even smaller final margin. This explains why an impressive
local idea may produce a very small asymptotic exponent in the whole algorithm.

\paragraph{Why strict slack is reserved.}
Let $G=\min_i g_i$. Fixed powers of $\log p$ can be absorbed if
$\kappa<G$, since $(\log p)^C=o(p^{G-\kappa})$ for fixed $C$.
The original conservative choice gives $G=2\kappa$. The final cost then fits
\[
 O(Vp^{1-\kappa})=O\bigl(n(\log n)^{1-\kappa}\bigr).
\]
This argument is eventual. A single fixed cutoff covers all the required
inequalities; below it, the machine uses schoolbook multiplication. Enlarging
the big-$O$ constant covers the finite set of smaller lengths. Establishing
the complete uniform implementation remains part of the proof obligation.

\paragraph{What later community refinements are trying to improve.}
Some refinements change a finite network or the price of its frame transitions;
others reuse workspace, change how stages are combined, or improve an outer
parameter choice. These operate at different layers. A new finite profile
must still be bound to its actual circuit and then passed through the applicable
recurrence and cost analysis. The original seven-margin diagram is a reading
aid for this frozen manuscript, not a claim that every later construction uses
the same accounting or that it displays the latest community witness.

\takeaway{Exact recovery needs a strict error bound; the exponent needs a
strict margin in every dominant cost. Both must refer to the same algorithm.}
\sourceanchor{08-assembly.tex, eq:fixed-parameters; eq:margin-list; eq:final-error}
